\par\bigskip\noindent\textbf{K101 supporting proof: The full grouped and alternating factor comparison}\par\nobreak
This local common modification retains both factor orders of EI.11--12,
SIT.4--5 and the ordinary O/T quotient. It supplies their simultaneous
corner morphism. The compact global attachment, its full fixed-face
ledger and its characteristic evaluation remain in progress. No
numerical return is assigned here.

\par

\medskip\noindent\textbf{FO.1. The two parents and their full coordinates}

\par\nobreak

Use the genuine point-line coordinate \(V\) and source graph
\(\sigma=\eta v\). On the selected overlap \(q\) is a unit. Its
saved nine-factor chart restricts to the grouped order
\[
 (\eta,\eta,\eta,\eta,v,v,v,v),\qquad UW=\eta^4v^4.             \tag{FO.1}
\]
The deleted \(q\) factor remains in the unit arrow below.
The ordinary associated-node resolution resolves each
\(X_kY_k=\eta v,\ V=(\eta v)^kX_k\) in order \((\eta,v)\).
Its four successive indices give the alternating order
\[
 (\eta,v,\eta,v,\eta,v,\eta,v).                              \tag{FO.2}
\]
Both modifications are retained.

In \(N_3=\mathbb Z^3\), ordered as \((\eta,v,V)\), write
\(A_j=(1,0,j)\) and \(C_j=(0,1,j)\). The alternating parents are
\[
 \langle C_j,A_{j+1},A_j\rangle,\quad
 \langle A_{j+1},C_{j+1},C_j\rangle \quad(0\le j\le3).         \tag{FO.3}
\]
The grouped parents are
\[
 \langle A_r,A_{r-1},C_0\rangle\ (1\le r\le4),\qquad
 \langle A_4,C_s,C_{s-1}\rangle\ (1\le s\le4).               \tag{FO.4}
\]
Their determinants have absolute value one. These rays follow
by substituting the ordered-factor coordinates
\(f_r=s_rt_r,\ U=s_r\prod_{\ell<r}f_\ell,\
W=t_r\prod_{\ell>r}f_\ell\).

On the slice \(\eta+v=1\), with \((x,y)=(\eta,V)\), both
supports are the rectangle
\[
 0\le x\le1,\qquad0\le y\le4.                               \tag{FO.5}
\]
The alternating walls are \(y=j,\ y=j+x\). The grouped
walls are \(y=rx,\ y=4x+s(1-x)\). Each of the two
partitions has disjoint interiors and covers this rectangle.
Its period is
\[
 S_3(\eta,v,V)=(\eta,v,V+4\eta+4v).                         \tag{FO.6}
\]

The full source/frame convention is retained. At selected
period index \(j\), its first line factor is
\[
 U_j=q^{-j}\psi\,(\eta v)^{-4j}V,\qquad
 \psi=f_{P^a}/r_a.                                        \tag{FO.7}
\]
Here \(\psi\) is the complete point-theta versus genuine
point-frame ratio of EI.77--79, evaluated on the actual
source graph. The period translation includes
\((\eta v)^{-4j}\); the unit \(q^{-j}\psi\) remains.
In the earlier chart notation the actual comparison is
\[
 X_{4j+r-1}=\psi^{-1}q^j s_r,\qquad
 Y_{4j+r-1}=\psi q^{-j}t_r.                                \tag{FO.8}
\]
Every integer index, Tate unit and Cartier ratio is retained.

\par

\medskip\noindent\textbf{FO.2. Exact common refinement on the \(\chi\)-unit overlap}

\par\nobreak

Intersect every triangle of FO.3 with every triangle of
FO.4. Exact rational clipping gives twenty-six
positive-area cells with twenty-three vertices. Their
interiors are disjoint, since each interior point lies
in one open triangle of each original partition. Their
union is FO.5 and their total area is four.

On each polygon edge insert every intersection vertex
lying on that edge. Triangulate the resulting convex
polygons by ears containing no other boundary vertex
on their diagonals or interiors. This gives thirty-four
cones. A rational vertex \((x,y)\) is lifted to the
primitive integral multiple of \((x,1-x,y)\). No lattice
cover is made. Every common boundary subdivision is
retained, so these triangulations agree on shared edges.

Two further stellar insertions suffice:
\[
 (2,1,7),\qquad(1,2,5).                                   \tag{FO.9}
\]
For each ray the two containing determinant-two cones
are replaced by four determinant-one cones. The final
common fan has thirty-eight regular maximal cones.
Every interior two-face occurs twice; exposed faces
lie on the four sides of FO.5. The lower face is
\(\langle A_0,C_0\rangle\), the upper face is precisely
\(S_3\langle A_0,C_0\rangle=\langle A_4,C_4\rangle\),
and both side chains remain unchanged.

The output saves every rational cell, triangle, primitive
ray, insertion and facet. This construction proves the
common-face property: intersections of the original
cones give a common fan; the polygon subdivisions
agree on their entire common edges; each later operation
subdivides every cone containing the same actual ray.

For either parent ray matrix \(P\) and a regular child
matrix \(N\), the actual parent coordinate map has
monomial matrix
\[
 P^{-1}N.                                                  \tag{FO.10}
\]
Every entry is a nonnegative integer. Both parent maps
are saved for every child and glue by the same torus
characters on common-face localizations.

In the grouped frame the complete coordinate formula is
\[
 y_i=(q^{-j}\psi)^{(P^{-1}e_V)_i}
       \prod_\nu z_\nu^{(P^{-1}N)_{i\nu}},                 \tag{FO.11}
\]
where the period-index part of \(P\) already includes
FO.6. Its inverse comparison retains \(q^j\psi^{-1}\).
The alternating parent already uses genuine \(V\).
All unit powers \(P^{-1}e_V\) are saved. These formulas
give the actual holomorphic maps, including the full
functions on the source graph.

Both maps are proper over their parent charts. Indeed,
a parent arc\textquotesingle s torus valuation lies in the parent cone
and hence in one of the finitely many child cones.
FO.10 extends its characters there. Two extensions
agree by the common-face localization. Applying the
same argument on the face of the generic orbit also
treats arcs whose generic point lies on a boundary
stratum. The finite glued chart map is separated and
satisfies this valuative criterion. Thus it is the
actual proper common modification.

\par

\medskip\noindent\textbf{FO.3. Extension through the actual fourth-root quotient}

\par\nobreak

In the selected \(q\)-unit neighbourhood \(\chi\) is a
unit, so the fourth-root action has no stabilizer
there. Toward the ordinary O/T corner use its actual
quotient cocharacter lattice
\[
 N_{\rm phys}=\mathbb Z^4+\mathbb Z(1,-1,0,0)/4
\]
with basis and dual invariant characters
\[
 ((1,-1,0,0)/4,e_\eta,e_v,e_V),\qquad
 (T_0,T_1,T_2,T_3)=(\chi^4,\chi\eta,v,V).                  \tag{FO.12}
\]
Embed a common ray as \((0,\eta,v,V)\) and adjoin
\[
 E=(4,1,0,0).                                              \tag{FO.13}
\]
Coning each of the thirty-eight common three-cones
with \(E\) gives thirty-eight four-cones of
determinant four. Since \(E\) lies outside their
common hyperplane, coning preserves every intersection
face. Each four-cone is contained in both raw parents
obtained by coning FO.3 and FO.4.
The second raw parent is the genuine-point-frame
extension of the grouped factor chart constructed
by this coning. On the \(\chi\)-unit overlap
FO.11 identifies it with the existing selected
chart. The function \(\psi\) is used on its
actual source-chart domain; no Tate frame is
required on an ordinary chart outside that domain.

The source period on this lattice is
\[
 S_4=
 \begin{pmatrix}
 1&0&0&0\\0&1&0&0\\0&0&1&0\\-1&4&4&1
 \end{pmatrix}.                                           \tag{FO.14}
\]
The period width of \(r\) is
\(-r_0+4r_1+4r_2\). The program inserts every period
image whose last coordinate lies between zero and
this width.

All regularization occurs in \(N_{\rm phys}\). The
exact residue group of each chosen cone is enumerated.
A nonzero point in its half-open fundamental
parallelepiped, divided by its positive common divisor,
gives the next primitive ray. Every containing cone is
subdivided, with all its period translates. The
calculation checks that each cone has at most one
nontrivial centre from a given period orbit and that
each positive coordinate of that centre is strictly
less than one. Every new determinant therefore strictly
decreases. Fifty-two periodic stellar steps give
two hundred thirty-eight cones of determinant \(+1\).

Every new quotient ray has positive first coordinate.
The entire restriction to the \(\chi\)-unit open
therefore stays unchanged: the complete three-face list
in the first-coordinate-zero hyperplane is exactly the
thirty-eight-cone fan of FO.2. There are four period
facets on each side, paired exactly by FO.14.

Here are the maps to both original coarse quotients.
For a raw parent \(P\), its actual four cover weights are
\[
 4P^{-1}(1,0,0,0)^t\pmod4.                                \tag{FO.15}
\]
List all nonnegative \(a\) with \(1\le|a|\le4\) and
weight zero. They generate its invariant polynomial
ring. To see this, order the factors in a larger
invariant monomial; among five successive partial
sums in \(\mathbb Z/4\), two agree, giving an invariant
block of at most four factors. Remove the block and
repeat. For every child \(N\), the full invariant
coordinate arrow is
\[
 x^a=\prod_\nu z_\nu^{(a^tP^{-1}N)_\nu}.                  \tag{FO.16}
\]
Every exponent is nonnegative integral, with complete
rows saved for both parents. Relations hold on the
dense quotient torus and thus on the whole chart.
The maps glue by their common characters and are
proper by the support and common-face criterion
already proved in FO.2.

This is an actual common quotient extension of the
two factor orders. The earlier fifty-two-cone
ordinary-only template and its audit remain preserved
as a separate local modification. The new common
template has 238 regular cones and both parent labels.

\par

\medskip\noindent\textbf{FO.4. Full source, Cartier, period and volume arrows}

\par\nobreak

The actual cover period is
\[
 (\eta,v,V)\longmapsto
 (\eta c_b,v/c_b,{\cal C}(\eta v)^4V),\qquad
 {\cal C}=\lambda H\mathcal B_{a,b}.                        \tag{FO.17}
\]
These are the full chord and genuine Hom functions
of EIR and EI at their actual intermediate source
points. On the invariant torus it becomes
\[
 (T_0,T_1,T_2,T_3)\longmapsto
 (T_0,c_bT_1,c_b^{-1}T_2,
      {\cal C}T_1^4T_2^4T_0^{-1}T_3).                     \tag{FO.18}
\]
For a regular child \(N\), the target is \(S_4N\).
The corresponding coordinates have identity monomial
map and the saved unit exponent vectors
\[
 (S_4N)^{-1}(0,1,-1,0)^t\quad\text{for }c_b,\qquad
 (S_4N)^{-1}(0,0,0,1)^t\quad\text{for }{\cal C}.             \tag{FO.19}
\]
Substitution recovers FO.18 exactly; the inverse
uses the carried inverse source arrow and units.
The grouped frame formula FO.11 is transported with
the target value of \(\psi\) and the actual quotient
\(q^{j+1}/q^j\). Both factors remain full functions.

The fourth-root source action and its point coefficient
give the quotient lattice FO.12. On the selected
\(\chi\)-unit charts its representative continues to
the next of the four selected branch neighbourhoods.
The exact source/theta functions carry FO.11 through
that arrow. No finite-root division section is used
on this ordinary overlap.

The physical circle is \((0,0,0,1)\), with chart weights
\[
 N^{-1}(0,0,0,1)^t.                                        \tag{FO.20}
\]
The original base functions are
\[
 \epsilon=T_1T_2,\quad q_{\rm new}=T_0,\quad
 \sigma^4=T_1^4T_2^4/T_0,\quad
 t_o=d(\epsilon)^{-3}\sigma^4.                             \tag{FO.21}
\]
Their nonnegative exponent rows are saved and satisfy
\(q_{\rm new}\sigma^4=\epsilon^4\) in every chart.

The actual original marked-family form is
\[
 2\Omega_o\wedge d\epsilon\wedge d\log V
 =8\eta^4v^3d^{-2}u_{O/T}\,
       d\chi\wedge d\eta\wedge dv\wedge d\log V.            \tag{FO.22}
\]
Since the last wedge is \(\epsilon/4\) times the
invariant logarithmic volume, the original form has
logarithmic character \((-1,5,4,0)\), with full unit
\(2d^{-2}u_{O/T}\) in that logarithmic frame.
On an oriented regular child its orders are
\[
 (-1,5,4,0)N-(1,1,1,1).                                   \tag{FO.23}
\]
The output retains these, the local-frame orders and
every physical weight. The full derivative square is
\[
 J_b=1+\eta\partial_\eta\log c_b-v\partial_v\log c_b,
 \qquad u_{\rm source}=u_{\rm target}c_bJ_b.                \tag{FO.24}
\]
Thus \(\eta^4v^3u\) transforms by the inverse source
Jacobian. In FO.11, \(d\log(q^{-j}\psi)\) is a
source/base one-form; its wedge with the three
source/base differentials is zero. This proves the
complete determinant-frame square including the
full \(q^j,\psi,c_b,\mathcal B\) factors.

There are two analytic O/T copies of this local
common template. The physical \(B\) exchanges them
and shifts the index by four; the fourth-root
generator fixes each centre separately. Period
faces therefore glue crosswise, retaining the actual
\(u_O/u_T,c_b\) factors. The separate two-torsion pair
has its own fixed-ideal \(\mu_2\) resolution. The
local cone count alone is not used for a global
fixed-component count.

\par

\medskip\noindent\textbf{FO.5. Reproducible output and remaining assembly}

\par\nobreak

The program
\path{code/part3/e6_selected_infinity_factor_overlay.py}
and its JSON output retain every parent, intersection
cell, primitive ray, cone, stellar step, period face,
invariant-generator arrow, physical weight, volume
order and full unit exponent. Its current run passes
21,224 exact assertions: 26 intersection cells,
38 regular common three-cones and 238 regular
common quotient four-cones. The full selected
extension of FO.6 has thirty-nine regular charts.

This supplies the missing factor-order interface.
The native I2 O/T graphs match by
\(\eta_e=\chi\eta,\ b_e=\chi\bar Y,\
\epsilon=\chi\sigma\), with the complete line
powers \(X_k=\chi^kA_k,\ Y_k=\chi^{-k-1}B_k\).
The original cap remains on its actual large-\(\sigma\)
domain, with word \(AB^{j_\gamma}\) and all retained
EIR period, Cartier and reference cochains. Those
domain and closing maps are unchanged by this
common modification.

\par

\medskip\noindent\textbf{FO.6. Extension to the full selected nine-factor chart}

\par\nobreak

The common \(q\)-unit interface extends through the
actual selected \(q=0\) chart. Use the integral
lattice ordered as \((\eta,v,q,V_\theta)\), and set
\[
 \iota(\eta,v,V)=(\eta,v,0,V),\qquad
 Q_0=(0,0,1,0),\quad Q_1=(0,0,1,1).                         \tag{FO.25}
\]
Cone each of the thirty-eight common three-cones
with \(Q_0\), and add
\[
 \langle\iota(A_4),\iota(C_4),Q_0,Q_1\rangle.                \tag{FO.26}
\]
These thirty-nine cones have determinant one.
They refine the actual nine-factor parent:
the first thirty-eight refine its eight
\(\eta\)- and \(v\)-factor cones; FO.26
is its unchanged final \(q\)-factor cone.
Their support is
\[
 \eta,v,q\ge0,\qquad0\le V_\theta\le4\eta+4v+q.              \tag{FO.27}
\]
The intersection of FO.26 with the lower region
is \(\langle\iota(A_4),\iota(C_4),Q_0\rangle\).
Thus every common face is retained. The period is
\[
 (\eta,v,q,V_\theta)\longmapsto
     (\eta,v,q,V_\theta+4\eta+4v+q).                        \tag{FO.28}
\]
It carries the lower face
\(\langle\iota(A_0),\iota(C_0),Q_0\rangle\)
onto the upper face
\(\langle\iota(A_4),\iota(C_4),Q_1\rangle\).
On the \(q\)-unit open the complete face set
with zero \(q\)-coordinate is the same
thirty-eight-cone fan. These assertions and
every nonnegative integral map to the
nine-factor parent are checked and saved.

Here \(V_\theta=\psi V\) is the actual
point-theta coordinate. The period-index
factor \((\sigma^4q)^{-j}\) is included
in FO.28. Consequently \(q^j\) is an actual
monomial period factor at \(q=0\); only on
the stated \(q\)-unit overlap does it enter
FO.7--11 as an invertible function.

The full ordinary and selected Hom factors satisfy
\[
 \beta(p)=\mathcal B(p)\frac{\psi(T_b p)}{\psi(p)},\qquad
 q=\kappa H,\qquad C_{\rm sel}=\lambda\beta/\kappa .
 \tag{FO.29}
\]
Substitute \(\psi=f_{P^a}/r_a\) into the
two complete Cartier ratios to verify this:
\[
 \mathcal B=\sigma^{-6}\bar B\,\frac{r_a(T_b p)}{r_a(p)},
 \qquad
 \beta=\sigma^{-6}\bar B\,
             \frac{f_{P^a}(T_b p)}{f_{P^a}(p)} .
\]
Thus the entire selected reset is
\[
 (\eta,v,q,V_\theta)\longmapsto
 (\eta c_b,v/c_b,q,C_{\rm sel}(\eta v)^4qV_\theta).
 \tag{FO.30}
\]
Its source- and line-unit vectors are the inverse
target matrix applied to \((1,-1,0,0)\) and
\((0,0,0,1)\), respectively. They are saved
for all thirty-nine charts. Equations
FO.29--30 make the ordinary/grouped/selected
square commute with the complete \(q^j,\psi\)
factors.

The source implicit coordinate \(\chi(\sigma,q)\)
has an extending nonzero derivative \(\chi_q\)
on the chosen selected neighbourhood. At its
centre \(H_\chi=-4\chi^3\) and
\(q_\chi=\kappa H_\chi\), so
\(\chi_q=-1/(4\kappa\chi^3)\), a unit.
The full derivative function enters the
original marked volume:
\[
 2\Omega_o\wedge d\epsilon\wedge d\log V_\theta
 =8\eta^4v^3d^{-2}u\,\chi_q\,
     d\eta\wedge dv\wedge dq\wedge d\log V_\theta.
 \tag{FO.31}
\]
Indeed \(d\sigma=v\,d\eta+\eta\,dv\), so the
\(\chi_\sigma d\sigma\) term vanishes in
FO.22; moving \(dq\) past \(d\eta,dv\)
makes two transpositions. Its logarithmic
character is \((5,4,1,0)\), and its actual
ray orders are
\((5,4,1,0)N-(1,1,1,1)\).
The full unit is \(8d^{-2}u\chi_q\).
All these orders are saved. FO.24 and the
full differential of FO.29 retain their
determinant-frame square.

All added modification centres have
\(\eta v=0\), hence \(\epsilon=0\).
The selected punctured family and its
physical closing are preserved through
the displayed inverse monomial and unit
maps. This adds the full selected chart
to the local common interface without
assigning a global characteristic number.

